\section{Validation of algorithms}
\label{sec:ebsd-hydrate-validation}

To validate the algorithms introduced in \zcref{sec:ebsd-hydrate,subsec:algo_hlm}, simulated datasets were created.
For this purpose, an artificial volume was randomly filled with \num{120000} spherical \ac{C3S} particles per \unit{\cubic\milli\meter}.
A \textsc{Gaussian} distribution of particle diameters was generated with diameters ranging from $d_\text{min, \ac{C3S}}=$\qty{0}{\micro\meter} to $d_\text{max, \ac{C3S}}=$\qty{5}{\micro\meter} was selected. 
The total volume was set to \qtyproduct{2000x2000}{\micro\meter\squared} with a height ($z$) of four times the maximum diameter of a hydrated \ac{C3S} grain.
The simulation allows the \ac{HLT} to be varied depending on the crystal orientation of the clinker grains.
To enable this, a random crystal orientation was assigned to each \ac{C3S} grain (based on the monoclinic unit cell \acs{ICSD} 94742 \cite{delatorre.2002}).
Two scenarios were then created for validation.
For the validation of the \ac{HLT} measurement algorithm, a constant spherical hydrate layer was used.
To validate the \ac{EBSD} correlation, the hydrate layer was modelled as an ellipsoid whose orientation varied according to the assigned crystal orientation.

Slicing this volume, generated a virtual \ac{BSE} image (\qty{0.2}{\micro\meter\per\px}), which was then be used as input for the methods to be validated.
Sections of such synthetic \ac{BSE} images are shown in \zcref{fig:c3s_synthetic_BSE}.

\begin{figure}[htb!]
    \centering

    \begin{subfigure}[t]{0.49\linewidth}
        \includegraphics[width=\textwidth,trim={0 0 0 1cm},clip]{hydrate-rim/ebsd/evaluation/synthetic_bse_slice_round.pdf}
        \caption{Spherical hydrate layer.}
    \end{subfigure}%
    \hfill%
    \begin{subfigure}[t]{0.49\linewidth}

        \includegraphics[width=\textwidth,trim={0 0 0 1cm},clip]{hydrate-rim/ebsd/evaluation/synthetic_bse_slice_eliptic.pdf}
        \caption{Ellipsoidal hydrate layer.}
    \end{subfigure}%
        
        \caption{
            Examples of synthetically generated \ac{BSE} images with a spherical \qty{2}{\micro\meter} hydrate layer (left) and an ellipsoidal hydrate layer with semi-axes of \qtyproduct{0.5x3.0x5}{\micro\meter} (right).
            The red lines depict the selected hydrate layer thicknesses.
            Black: porosity; dark grey: \ac{CSH}; light grey: \ac{C3S}.
        }
    \label{fig:c3s_synthetic_BSE}
\end{figure}

Applying the \ac{HLT} measurement to the synthetic \ac{BSE} images results in \zcref{fig:c3s_synthetic_BSE_HLT}.
The diagrams consist of a graph indicating frequency of a measured \ac{HLT} as a function of the grain size (blue-green).
This plot is accompanied by a frequency plot of the grain size diameter (grey) and a grain-size-agnostic distribution plot of the \ac{HLT} frequency (orange).
The set values for the respective \ac{HLT} are marked with red lines.

\begin{figure}[htb!]
    \centering

    \includegraphics[width=\textwidth]{hydrate-rim/ebsd/evaluation/2D_cir_vs_elliptic.pdf}
        
    \caption{
        Normalised distributions of the \ac{HLT} of simulated hydrated \ac{C3S} depending on the particle diameter; 
        left: spherical hydrate layer (\qty{3}{\micro\meter}); right: ellipsoidal hydrate layer (\qtyproduct{0.5x3.0x5}{\micro\meter}).
        The histogram (grey) on top illustrates the logarithmic frequency for each diameter bin in the diagram below.
        The histogram (orange) on the right illustrates the frequency of \ac{HLT} measurements as the sum of the main diagram.
    }
    \label{fig:c3s_synthetic_BSE_HLT}
\end{figure}

The diagrams are distinctly different.
The spherical hydrate layer results in a fairly constant \ac{HLT} with a maximum at \qty{3.3}{\micro\meter}.
This value is \qty{10}{\percent} higher than the simulated layer thickness.
While one reason for this overestimation is the sectioning effect (see \zcref{subsec:sectioning-effect}), another reason is the merging of the hydrate layers of multiple grains, causing a larger \ac{HLT} for some measurements.
The simulated elliptical hydrate layer shows a much broader distribution of the \ac{HLT}.
The maximum value of the \ac{HLT} was measured at \qty{2.9}{\micro\meter}, but the majority of measurements are distributed between the set limits for the ellipsoid.

For both measurements, the data density below a grain diameter of \qty{4}{\micro\meter} was low.
However, these smaller grains seem to have a tendency to larger \ac{HLT}s, which can also be explained with the sectioning effect.
In real specimens, many more smaller grains are expected than in this simulation, changing the appearance of the distribution for smaller grains.
Finally, the 2D plot and the grain diameter histogram show regular vertical artefacts, which are probably caused by rounding errors, but do not affect the evaluation and are less pronounced in real-world data.

Using the assigned crystal orientations, a simulated \ac{EBSD} dataset was generated with a resolution of \qty{2}{\micro\meter\per\px} in the form of a channel text file.
The resolution of the simulated \ac{BSE} image only had to be adjusted to the \ac{EBSD} dataset, since the datasets were already aligned to each other.
After processing the data as explained in \zcref{sec:ebsd-hydrate}, the \ac{HLT} measurements can be projected onto a unit sphere representing the crystal orientation of \ac{C3S} (see \zcref{fig:ebsd_eval_sphere}).

\begin{figure}[htb!]
	\centering
    
    \begin{subfigure}[t]{0.49\linewidth}
        \begin{tikzpicture}
            \begin{axis}[
                width=0.66\textwidth,
                height=0.66\textwidth,
                scale only axis,
                enlargelimits=false,
                xmin=-1.1, xmax=1.1,
                ymin=-1.1, ymax=1.1,
                axis on top,
                colormap name=magma,
                colorbar,
                colorbar style={
                    ylabel={\ac{HLT} in \unit{\micro\meter}},
                },
                point meta min=0,
                point meta max=12.1,
                ticklabel style={font=\small},
                xtick={-1,0,1},
                ytick={-1,0,1},
            ]
                \addplot [draw=none] coordinates {(-1,-1) (1,1)};
                \node[inner sep=0pt, anchor=south west] at (axis cs:-1,-1) {
                    \ifdefined\PRINTABLE
                        \includegraphics[width=0.6\textwidth]{hydrate-rim/ebsd/evaluation/ani_scatter/scatter00.jpg}
                    \else
                        \animategraphics[loop,autoplay,width=0.6\textwidth]{12}{hydrate-rim/ebsd/evaluation/ani_scatter/scatter}{00}{35}
                    \fi
                };
            \end{axis}
        \end{tikzpicture}
        \caption{
            \ac{HLT} scatter plot on a unit sphere..
            \label{fig:ebsd_eval_sphere_a}
        }
    \end{subfigure}%
    \hfill%
    \begin{subfigure}[t]{0.49\linewidth}
        
        \begin{tikzpicture}
            \begin{axis}[
                width=0.66\textwidth,
                height=0.66\textwidth,
                scale only axis,
                enlargelimits=false,
                xmin=-1.1, xmax=1.1,
                ymin=-1.1, ymax=1.1,
                axis on top,
                colormap name=magma,
                colorbar,
                colorbar style={
                    ylabel={\ac{HLT} in \unit{\micro\meter}},
                },
                point meta min=0,
                point meta max=4.07,
                xtick={-1,0,1},
                ytick={-1,0,1},
                ticklabel style={font=\small},
            ]
                \addplot [draw=none] coordinates {(-1,-1) (1,1)};
                \node[inner sep=0pt, anchor=south west] at (axis cs:-1,-1) {
                    
                    \ifdefined\PRINTABLE
                        \includegraphics[width=0.6\textwidth]{hydrate-rim/ebsd/evaluation/ani_healpix/healpix04.jpg}
                    \else
                        \animategraphics[loop,autoplay,width=0.6\textwidth]{12}{hydrate-rim/ebsd/evaluation/ani_healpix/healpix}{35}{00}
                    \fi
                    
                };
            \end{axis}
        \end{tikzpicture}
        \caption{
            Mean \ac{HLT} in a \textsc{HEALPix} grid on a unit sphere.
            \label{fig:ebsd_eval_sphere_b}
        }
    \end{subfigure}
    
	\caption{
        Unit sphere indicating \ac{HLT} measurements corresponding to the crystal orientation of a simulated \ac{C3S} dataset with an elliptical hydrate layer (\qtyproduct{0.5x3.0x5}{\micro\meter}) in a scatter plot (a) and in a \ac{HEALPix} grid (b).
        \label{fig:ebsd_eval_sphere}
    }
\end{figure}

Since the measurement density is not very high, the unit sphere was divided into a so-called \ac{HEALPix} grid \cite{gorski.2005}.
In a \ac{HEALPix} grid, the surface of the sphere is divided into sections of equal area.
This avoids a statistical distortion, especially at the poles of the unit sphere.
The mean \ac{HLT} within each \ac{HEALPix} grid cell was calculated and each cell was coloured accordingly (see \zcref{fig:ebsd_eval_sphere_b}).

If there is a dependence between the crystal orientation and the \ac{HLT}, it should become visible as a concentration of longer or shorter \ac{HLT}s in certain areas of the unit sphere.
Additionally, for an ideal preparation a homogeneous distribution of measurement counts is expected around the unit sphere.
While it becomes apparent for the elliptical hydration rim simulation, that the \ac{HLT} distribution is not uniform in \zcref{fig:ebsd_eval_sphere_b}, a \textsc{Mollweide} projection of the whole sphere surface is shown in \zcref{fig:ebsd_eval_sphere_projection}.

\begin{figure}[htb!]
	\centering
    \begin{subfigure}[t]{0.49\linewidth}
        
        \begin{tikzpicture}
            \begin{axis}[
                width=.8\textwidth,
                height=.45\textwidth,
                scale only axis,
                enlargelimits=false,
                colormap name=magma,
                colorbar,
                colorbar style={
                    ylabel={Average \ac{HLT} in \unit{\micro\meter}},
                    font=\small
                },
                point meta min=0,
                point meta max=4.07,
                xtick={-1,0,1},
                ytick={-1,0,1},
                ticklabel style={font=\footnotesize},
                hide axis,
            ]
                \addplot [draw=none] coordinates {(-1,-1) (1,1)};
                \node[inner sep=0pt, anchor=south west] at (axis cs:-1,-1) {
                    
                \includegraphics[width=.8\textwidth]{hydrate-rim/ebsd/evaluation/ellipsoid_mollweide_len.png}
                    
                };
            \end{axis}
        \end{tikzpicture}
        
        \caption{
            Mean \ac{HLT}.
            \label{fig:ebsd_eval_sphere_projection_b}
        }

    \end{subfigure}
    \hfill
    \begin{subfigure}[t]{0.49\linewidth}
        \begin{tikzpicture}
            \begin{axis}[
                width=.8\textwidth,
                height=.45\textwidth,
                scale only axis,
                enlargelimits=false,
                hide axis,
                colormap name=ocean_r,
                colorbar,
                colorbar style={
                    ylabel={$log_{10}(n+1)$},
                    font=\small
                },
                point meta min=0,
                point meta max=2.25,
                xtick={-1,0,1},
                ytick={-1,0,1},
                ticklabel style={font=\footnotesize},
            ]
                \addplot [draw=none] coordinates {(-1,-1) (1,1)};
                \node[inner sep=0pt, anchor=south west] at (axis cs:-1,-1) {
                    
                \includegraphics[width=.8\textwidth]{hydrate-rim/ebsd/evaluation/ellipsoid_mollweide_cnt.png}
                    
                };
            \end{axis}
        \end{tikzpicture}
        
        \caption{
            Data density plot.
            \label{fig:ebsd_eval_sphere_projection_a}
        }
    \end{subfigure}
    
	\caption{
        The same \ac{HEALPix} grid as in \zcref{fig:ebsd_eval_sphere_b} is shown in a \textsc{Mollweide} projection (a), together with the corresponding data density plot (b). 
        The lines ({\color{red}red}) are a projection of the monoclinic structure of \ac{C3S} as derived from \zcref{fig:c3s_unit_cell}. % HEALPix Mollweide
        \label{fig:ebsd_eval_sphere_projection}
    }
\end{figure}

This visualisation underlines that, for the elliptical hydration rim simulation, both the inhomogeneous but symmetric \ac{HLT} distribution and the homogeneous data density distribution are present as expected.
The \textsc{Mollweide} projection from \zcref{fig:ebsd_eval_sphere_projection_b} can also be projected back to the \ac{C3S} unit cell as illustrated in \zcref{fig:c3s_eval_backprojection}.
This figure shows the unit cell in 3D and unfolded as a net, and demonstrates which crystallographic surface was set to grow the largest \ac{HLT} value in the simulation.

\begin{figure}[htb!]
    \centering

    \begin{tikzpicture}
        \begin{axis}[
            width=.87\textwidth,
            height=.27\textwidth,
            scale only axis,
            enlargelimits=false,
            hide axis,
            colormap name=magma,
            colorbar,
            colorbar style={
                ylabel={Average \ac{HLT} in \unit{\micro\meter}},
            },
            point meta min=0,
            point meta max=4.07,
            xtick={-1,0,1},
            ytick={-1,0,1},
            ticklabel style={font=\small},
        ]
            \addplot [draw=none] coordinates {(-1,-1) (1,1)};
            \node[inner sep=0pt, anchor=south west] at (axis cs:-1,-1) {
                
            \includegraphics[width=.9\textwidth]{hydrate-rim/ebsd/evaluation/healpix_meanlen_cube_net_2d+3D.pdf}
                
            };
        \end{axis}
    \end{tikzpicture}
        
    \caption{
        Projection of the unit sphere from \zcref{fig:ebsd_eval_sphere_projection} to the unit cell of \ac{C3S}, MIII polymorph.
        Left: 3D visualisation of the unit cell with the mean \ac{HLT} values on each surface.
        Right: Net of the unit cell with the projected HEALPix grid.
    }
    \label{fig:c3s_eval_backprojection}
\end{figure}

The same experiment was done for the spherical hydrate layer.
As expected, both the \ac{HLT} and the measurement density were distributed homogeneously on the unit sphere.